% ID: FIS-CIR-SER-001
% Titolo: Due resistenze collegate in serie e in parallelo
% Disciplina: Fisica
% Area: Circuiti elettrici
% Argomento: Circuiti serie parallelo
% Sottoargomento: Resistenza equivalente e corrente totale
% Classe: Prima istituto professionale
% Difficolta: 3
% Tipo: problema_numerico
% Risultato: R_s = 18 Ohm; I_s circa 0,67 A; R_p = 4 Ohm; I_p = 3,0 A
% Tempo_stimato: 9 min
% Competenze: calcolo della resistenza equivalente, applicazione della legge di Ohm, confronto tra configurazioni
% Prerequisiti: resistenza elettrica, legge di Ohm, circuiti semplici
% Tag: fisica, circuiti_elettrici, serie, parallelo, legge_ohm
% Fonte: esempio_originale
% Autore: OpenAI Codex
% Licenza: CC-BY-SA-4.0

\begin{esercizio}
Due resistenze, \(R_1=6\,\Omega\) e \(R_2=12\,\Omega\), vengono collegate
a un generatore da \(12\,\mathrm{V}\).

Calcola la resistenza equivalente e la corrente totale se le due resistenze
sono collegate prima in serie e poi in parallelo.
\end{esercizio}

\begin{soluzione}
Nel collegamento in serie la resistenza equivalente e la somma delle due
resistenze:
\[
R_s=R_1+R_2=6\,\Omega+12\,\Omega=18\,\Omega.
\]

La corrente totale si trova con la legge di Ohm:
\[
I_s=\frac{V}{R_s}=\frac{12}{18}\,\mathrm{A}\simeq 0{,}67\,\mathrm{A}.
\]

Nel collegamento in parallelo vale:
\[
\frac{1}{R_p}=\frac{1}{R_1}+\frac{1}{R_2}.
\]

Quindi:
\[
\frac{1}{R_p}=\frac{1}{6}+\frac{1}{12}
=\frac{2}{12}+\frac{1}{12}
=\frac{3}{12}
=\frac{1}{4}.
\]

\[
R_p=4\,\Omega.
\]

La corrente totale in parallelo e:
\[
I_p=\frac{V}{R_p}=\frac{12}{4}=3{,}0\,\mathrm{A}.
\]

A parita di tensione, il collegamento in parallelo ha resistenza equivalente
minore e quindi corrente totale maggiore.
\end{soluzione}
